\chapter{TÓM LƯỢC}
\indent\indent Trong thời đại số hóa, các văn bản kỹ thuật số như hóa đơn tài chính, hợp đồng pháp lý, chứng chỉ học thuật và căn cước công dân đóng vai trò nền tảng trong các hoạt động giao dịch xã hội. Tuy nhiên, sự phát triển của các công cụ chỉnh sửa đồ họa đã khiến tình trạng làm giả văn bản (tampered text) trở nên ngày càng tinh vi, đe dọa trực tiếp đến tính toàn vẹn của các hệ thống xác thực điện tử. Các phương pháp pháp y truyền thống thường khó nhận diện các chỉnh sửa vi mô sau nhiều chu kỳ nén, trong khi các mạng học sâu hiện hành đa phần chỉ giải quyết bài toán đơn nhiệm (hoặc chỉ phân loại nhị phân thật/giả, hoặc chỉ phân đoạn pixel) với số lượng tham số khổng lồ, khó triển khai trên các thiết bị giới hạn tài nguyên.

Nhằm giải quyết triệt để các thách thức trên, đề tài đề xuất kiến trúc học sâu đa nhiệm gọn nhẹ mang tên \textbf{LightweightDTD} với tổng số tham số chỉ khoảng \textbf{4.15 triệu tham số (4,155,796 tham số)}. Hệ thống kết hợp hài hòa giữa hai luồng thông tin bổ trợ: (1) \textbf{Nhánh không gian (Spatial Branch)} sử dụng backbone EfficientNet-B0 được tối ưu hóa theo phương pháp Co giãn phức hợp để trích xuất đặc trưng hình ảnh đa tầng; và (2) \textbf{Nhánh tần số (Frequency Branch)} ứng dụng biến đổi cosin rời rạc theo khối \textbf{BlockDCT $8 \times 8$} trực chuẩn nhằm bắt trọn các bất thường trong ma trận nén JPEG và dấu vết nội suy vi mô. Hai nguồn đặc trưng được hòa trộn qua mô-đun chú ý không gian -- kênh đồng thời \textbf{scSE (Concurrent Spatial and Channel Squeeze-and-Excitation)} và đưa vào bộ giải mã lặp đa góc nhìn \textbf{LiteMID (Lightweight Multi-view Iterative Decoder)} để tích hợp ngữ cảnh đa tỉ lệ ($H/4$ đến $H/32$). Mô hình thực hiện cơ chế học đa nhiệm (Multi-task Learning) hỗ trợ song song hai đầu ra: phân đoạn vị trí làm giả ở mức pixel ($224 \times 224$) và phân loại đa nhãn 6 đặc trưng pháp y vi mô.

Thực nghiệm trên tập dữ liệu chuẩn \textbf{Minh-TFR TestSet} (gồm 2,000 mẫu kiểm thử độc lập) chứng minh tính vượt trội của mô hình: ở nhiệm vụ phân đoạn pixel, mô hình đạt \textbf{Độ chính xác pixel 99.79\%}, \textbf{Precision (Micro) đạt 90.09\%}, \textbf{Recall (Micro) đạt 77.93\%}, \textbf{IoU (Micro) đạt 71.78\%} và \textbf{Macro Dice đạt 62.18\%}. Ở nhiệm vụ phân loại đa nhãn đặc trưng làm giả, mô hình đạt \textbf{Micro F1 77.49\%}, \textbf{Macro F1 58.60\%}, \textbf{mAP (PR-AUC) đạt 61.02\%} và \textbf{Top-3 Class Hit Rate đạt tới 98.45\%}. Kết quả khẳng định khả năng phát hiện chính xác, định vị rõ nét vùng văn bản giả mạo với chi phí tính toán cực kỳ tiết kiệm, mở ra triển vọng ứng dụng lớn trong thực tế giám định số.\vspace{0.3cm}

\textbf{Từ khóa:} Giám định tài liệu số, Phát hiện chữ bị làm giả, LightweightDTD, EfficientNet-B0, BlockDCT, scSE Module, LiteMID, Học đa nhiệm, Multi-label F1.
